% ======================================================================
% INTRODUCTION CHAPTER
% ----------------------------------------------------------------------
% I did not add a \begin{chapterabstractpage} here because this is the
% introduction chapter, unless you do want/need that, then by all means
% add it.% 
% ======================================================================

\section{Introduction}
Taken from the Maastricht University Doctoral Regulations (Chapter 4, Article 12.2):  
\begin{quote}
If the thesis consists of one single scientific treatise, it must have the following components:
\begin{itemize}
  \item an introduction indicating the position of the research compared with other related research in a national or international context;
  \item the scientific treatise regarding a certain subject;
  \item a general discussion consisting of a reasoned representation of the doctoral candidate's own position with regard to the main theme, or the major themes of the thesis.
\end{itemize}
If the thesis consists of several individual scientific treatises, they must have sufficient cohesion. The cohesion between treatises must be explained by the doctoral candidate in the introduction and/or the general conclusion. At least one scientific treatise must be an article that demonstrably is in the review process of a scientific journal.
\end{quote}

(See: \href{https://umployee.maastrichtuniversity.nl/groep-media/um-policy-hub/doctoral-regulations-2023-including-annexes}{UM Policy Hub — Doctoral Regulations 2023, including annexes})

\bigskip

\noindent
So, this chapter should provide a clear introduction to the doctoral research. It should outline the topic, situate it within the broader field (preferably both nationally and internationally), explain its scientific and societal relevance, and describe how the individual studies (aka. treatises) contribute to the overall research question. \textbf{Make sure to replace this placeholder text with your own introduction} once you start writing your thesis.